\documentclass[12pt,a4paper]{article}
\usepackage{amsmath,amssymb,amsthm,booktabs}
\usepackage{hyperref}
\usepackage{geometry}
\geometry{margin=2.5cm}
\usepackage{enumitem}

\newtheorem{theorem}{Theorem}[section]
\newtheorem{lemma}[theorem]{Lemma}
\newtheorem{corollary}[theorem]{Corollary}
\newtheorem{proposition}[theorem]{Proposition}
\theoremstyle{definition}
\newtheorem{definition}[theorem]{Definition}
\theoremstyle{remark}
\newtheorem{remark}[theorem]{Remark}

\title{The MiniBooNE/LSND Anomaly and Sterile Neutrino Exclusion\\
from $D = 3$}

\author{Jonathan Washburn\\
Recognition Science Research Institute\\
Austin, Texas\\
\texttt{washburn.jonathan@gmail.com}}

\date{March 2026}

\begin{document}
\maketitle

\begin{abstract}
We analyze the MiniBooNE/LSND anomaly within Recognition Science (RS).
Both experiments report excess $\bar{\nu}_e$ appearance in $\nu_\mu$
beams, consistent with neutrino oscillations at $\Delta m^2 \sim
1~\mathrm{eV}^2$---a mass splitting incompatible with the standard
three-neutrino framework.  The conventional explanation invokes a
sterile neutrino with $m_4 \sim 1$~eV.  RS \emph{excludes} this
interpretation: the spatial dimension $D = 3$ forces exactly three
neutrino generations via the face-pair count of the recognition cube,
$N_\nu = \mathrm{face\_pairs}(3) = 3$, leaving no structural slot for a
fourth mass eigenstate.  We further propose that the $\varphi$-ladder
spacing of the three active neutrino masses
($m_i \propto \varphi^{n_i}$) produces $L/E$-dependent interference
modulations that can, at specific baselines, mimic a sterile oscillation
signal.  This mechanism is a hypothesis, not a proved theorem; it is
offered as a falsifiable alternative to sterile neutrino models.  The
MicroBooNE experiment (2021) has demonstrated that a significant fraction
of the MiniBooNE excess can be explained by standard-model backgrounds
(radiative $\Delta \to N\gamma$ decays); the RS analysis is consistent
with this finding.  The structural theorem ($N_\nu = 3$) follows from the RS forcing
chain; the $\varphi$-aliasing hypothesis awaits full quantitative
treatment.  The Short-Baseline Neutrino (SBN) program at
Fermilab---SBND, MicroBooNE, and ICARUS---provides the decisive
experimental test.
\end{abstract}

%% ============================================================
\section{Recognition Science Framework}
\label{sec:framework}
%% ============================================================

The structural exclusion of sterile neutrinos follows from the RS
forcing chain, which originates in the cost axioms.

\begin{definition}[RS Cost Functional]
For $x > 0$: $J(x) = \tfrac{1}{2}(x + x^{-1}) - 1$.
\end{definition}

\noindent\textbf{Axiom A1 (Normalization):} $J(1) = 0$.\\
\textbf{Axiom A2 (Recognition Composition Law, RCL):}
$J(xy)+J(x/y)=2J(x)J(y)+2J(x)+2J(y)$ for all $x,y>0$.\\
\textbf{Axiom A3 (Calibration):} $J''_{\log}(0) = 1$, where
$J_{\log}(t) := J(e^t)$.

\begin{theorem}[Cost Uniqueness]
\label{thm:unique}
The continuous solution to \textup{(A1)--(A3)} is unique:
$J(x) = \tfrac{1}{2}(x + x^{-1}) - 1$.
\end{theorem}

\begin{proof}[Proof sketch]
Set $H(t) = J_{\log}(t)+1$.  Axiom A2 transforms into the d'Alembert
equation $H(s+t)+H(s-t)=2H(s)H(t)$.  By the Acz\'el classification
theorem~\cite{Aczel1966}, the unique continuous solution with $H(0)=1$
and $H''(0)=1$ is $H(t)=\cosh t$, giving
$J(x)=\tfrac{1}{2}(x+x^{-1})-1$.
\end{proof}

\noindent The recognition cube $Q_D$ has spatial dimension $D=3$ forced
by $2^D=8$ (the 8-tick epoch) and three independent dimension-forcing
arguments.  Its face-pair count $\mathrm{face\_pairs}(3)=3$ is the
structural origin of exactly three active neutrino generations.  The
$\varphi$-ladder ($\varphi=(1+\sqrt{5})/2$ forced by RCL self-similarity)
assigns masses $m_i \propto \varphi^{n_i}$ to the three mass eigenstates.
No fourth mass eigenstate is admitted by the RS cube structure.

%% ============================================================
\section{Introduction}
\label{sec:intro}
%% ============================================================

The three-neutrino oscillation framework successfully describes solar,
atmospheric, and long-baseline reactor data with two independent mass
splittings: $\Delta m^2_{21} \approx 7.5 \times 10^{-5}~\mathrm{eV}^2$
(solar) and $|\Delta m^2_{32}| \approx 2.5 \times 10^{-3}~\mathrm{eV}^2$
(atmospheric)~\cite{PDG}.

Two experiments report anomalous $\bar{\nu}_e$ appearance that does
not fit this framework:
\begin{itemize}[nosep]
  \item \textbf{LSND} (1996--2001): excess $\bar{\nu}_e$ in a
  $\bar{\nu}_\mu$ beam from $\mu^+$ decay at rest; $L \approx 30$~m,
  $E \sim 20$--$200$~MeV; $\Delta m^2 \sim
  1~\mathrm{eV}^2$~\cite{LSND}.
  \item \textbf{MiniBooNE} (2018, 2020): excess low-energy $\nu_e$-like
  events in a $\nu_\mu$ beam at Fermilab; $L \approx 500$~m,
  $E \sim 200$--$3000$~MeV; combined significance ${\sim}\,4.8\sigma$
  after 12.84$\times 10^{20}$ protons on target~\cite{MiniBooNE2020}.
\end{itemize}

The combined LSND+MiniBooNE significance exceeds $6\sigma$.  The
canonical new-physics interpretation is a sterile neutrino with
$m_4 \sim 1$~eV and mixing $|U_{e4}|^2 \sim 10^{-3}$, forming a
``3+1'' oscillation model.  This is in tension with:
\begin{itemize}[nosep]
  \item Muon-neutrino disappearance bounds (IceCube, MINOS+).
  \item Cosmological limits on $N_{\mathrm{eff}}$ and the sum of
  neutrino masses.
  \item MicroBooNE's identification of radiative backgrounds in the
  same energy region~\cite{MicroBooNE2021}.
\end{itemize}

The status of the anomaly is therefore \emph{ambiguous}: it may reflect
an uncontrolled background, a systematic error, or genuine new physics
incompatible with the standard three-neutrino picture.  The RS
analysis~\cite{RS_Axioms} below offers a structural argument against
sterile neutrinos and a
qualitative model for residual excess events.

%% ============================================================
\section{The Structural Argument: $N_\nu = 3$ from $D = 3$}
\label{sec:structure}
%% ============================================================

\begin{theorem}[Exactly Three Neutrino Generations]
\label{thm:3nu}
In Recognition Science, the spatial dimension $D = 3$ forces exactly
three active neutrino generations:
\begin{equation}
  N_{\nu} = \mathrm{face\_pairs}(D) = D = 3.
\end{equation}
A fourth (sterile) neutrino is excluded by the spinor structure of the
$D$-dimensional recognition cube.
\end{theorem}

\begin{proof}
The recognition cube $Q_D$ has $2D$ faces organized into $D$ orthogonal
face pairs (top/bottom, front/back, left/right for $D = 3$).  Each
face pair generates one independent chiral spinor mode; these are the
three active neutrino flavors $\nu_e, \nu_\mu, \nu_\tau$.  With $D = 3$:
$\mathrm{face\_pairs}(3) = 3$.  There is no fourth face pair; hence
there is no structural slot for a fourth light neutrino species.
\end{proof}

\begin{corollary}
$\mathrm{face\_pairs}(3) \neq 4$: a sterile neutrino at $m_4 \sim 1$~eV
is excluded by RS.  The correct prediction from RS is $N_\nu = 3$
active flavors with masses on the $\varphi$-ladder.
\end{corollary}

\begin{remark}
The dimension $D = 3$ is itself forced in RS by three independent
arguments derived from the RS forcing chain~\cite{RS_Axioms}: (i)
Alexander duality requires integer linking invariants only in $D = 3$;
(ii) Kepler stability of near-circular orbits requires $D = 3$;
(iii) the minimum dyadic-to-sexagesimal synchronization (8-tick and
45-tick cycles) achieves $\mathrm{lcm}(8, 45) = 360$ exactly at $D = 3$.
\end{remark}

%% ============================================================
\section{The $\varphi$-Aliasing Hypothesis}
\label{sec:aliasing}
%% ============================================================

The exclusion of sterile neutrinos raises the question: what produces
the LSND/MiniBooNE excess?  We propose, as a hypothesis to be tested
quantitatively, the following mechanism.

\begin{proposition}[$\varphi$-Ladder Neutrino Masses]
\label{prop:nu-masses}
In RS, the three active neutrino masses sit on the $\varphi$-ladder:
\begin{equation}
  m_i \approx m_0 \, \varphi^{n_i}, \quad n_1 < n_2 < n_3,
\end{equation}
where $m_0$ is a reference mass set by the coherence energy scale and
$\varphi = (1+\sqrt{5})/2$.
\end{proposition}

\begin{definition}[$\varphi$-Aliasing]
\label{def:aliasing}
At baseline $L$ and energy $E$, the two-flavor oscillation probability
$P(\nu_\mu \to \nu_e)$ receives contributions from all three mass
splittings $\Delta m^2_{ij}$.  For $\varphi$-spaced masses, higher
rungs give $\Delta m^2_{3i} \sim \varphi^{2(n_3 - n_i)} \Delta
m^2_{21}$.  At baselines $L/E$ tuned to a high rung splitting, the
summed oscillation probability can exhibit a beat frequency that
is numerically consistent with a \emph{new} splitting $\Delta
m^2_{\mathrm{eff}} \sim 1~\mathrm{eV}^2$, without introducing a fourth
eigenstate.
\end{definition}

\begin{remark}
This mechanism is qualitatively analogous to aliasing in signal
processing: the discrete $\varphi$-ladder spacing can create apparent
``extra'' oscillation frequencies in a finite baseline window.  However,
the analogy requires a quantitative calculation to determine whether the
$\varphi$-spaced active masses can actually produce an apparent
$\Delta m^2 \sim 1~\mathrm{eV}^2$ at the LSND and MiniBooNE baselines.
This calculation is \emph{not} completed in the present paper and
represents an open question.  The hypothesis is stated here as a
target for future work.
\end{remark}

%% ============================================================
\section{MicroBooNE and Background Explanations}
\label{sec:microboone}
%% ============================================================

The MicroBooNE experiment~\cite{MicroBooNE2021} operated in the same
neutrino beam as MiniBooNE with a liquid-argon time-projection chamber
that allows event-by-event particle identification.  Key findings:

\begin{itemize}
  \item MicroBooNE found no excess of charged-current quasi-elastic
  $\nu_e$ events consistent with a simple sterile neutrino oscillation.
  \item Radiative $\Delta^+ \to p + \gamma$ decays, in which the photon
  mimics an electron in the MiniBooNE Cherenkov detector, account for
  a substantial fraction of the observed excess in the MiniBooNE signal
  region.
  \item A residual excess may remain in some MicroBooNE channels, but
  its nature and significance continue to be analyzed.
\end{itemize}

\begin{theorem}[RS Consistency with Background Explanations]
\label{thm:bg-ok}
If the MiniBooNE excess is entirely due to mis-identified
$\Delta \to N\gamma$ and $\pi^0$ backgrounds, the three-generation RS
prediction is confirmed: no fourth neutrino eigenstate exists.
\end{theorem}

\begin{proof}
The background explanation requires only standard-model processes and
the three active neutrinos.  This is entirely consistent with
$N_\nu = 3$ (Theorem~\ref{thm:3nu}).
\end{proof}

%% ============================================================
\section{Comparison with Other Explanations}
\label{sec:compare}
%% ============================================================

\begin{center}
\renewcommand{\arraystretch}{1.3}
\begin{tabular}{p{3.2cm}p{3.8cm}p{5.2cm}}
\toprule
\textbf{Explanation} & \textbf{Prediction} & \textbf{Status} \\
\midrule
Sterile neutrino (3+1) & $\nu_4$ with $m_4 \sim 1$~eV & Tension with
IceCube, MINOS+, $N_{\rm eff}$; excluded by RS ($N_\nu = 3$) \\
SM backgrounds ($\Delta$, $\pi^0$) & No fourth mass eigenstate & 
Supported by MicroBooNE; consistent with RS \\
Lorentz violation & CPT-odd oscillations & No evidence; not predicted by RS \\
RS $\varphi$-aliasing & Three-active-$\nu$ beat frequency & 
Hypothesis; quantitative calculation needed \\
\bottomrule
\end{tabular}
\end{center}

%% ============================================================
\section{Predictions and Falsifiers}
\label{sec:predict}
%% ============================================================

\begin{enumerate}
  \item \textbf{No sterile neutrino discovery.}  The SBN program
  (SBND, MicroBooNE, ICARUS) will not find evidence for a
  fourth neutrino mass eigenstate.  Any residual excess will be
  consistent with backgrounds or the $\varphi$-aliasing mechanism,
  not a new oscillation frequency.  \emph{Falsifier}: a confirmed
  oscillation minimum at the predicted $\Delta m^2_{14}$ in a
  disappearance channel.

  \item \textbf{$N_{\mathrm{eff}} = 3.044$.}  The effective number of
  thermalized neutrino species remains at the standard three-generation
  value.  A measurement of $N_{\mathrm{eff}} > 3.1$ at $>5\sigma$
  would challenge RS.

  \item \textbf{$\varphi$-aliasing test.}  If the $\varphi$-aliasing
  hypothesis is developed quantitatively, it predicts a specific $L/E$
  dependence for any residual excess---one consistent with an
  interference pattern of three $\varphi$-spaced oscillation
  frequencies, not a simple two-frequency pattern.  High-statistics SBN
  data can distinguish these.

  \item \textbf{Background-only explanation.}  If MicroBooNE's full
  analysis and the SBND data find that all excess events are accounted
  for by known backgrounds, the RS three-generation prediction stands.
\end{enumerate}

%% ============================================================
\section{Conclusion}
\label{sec:conclusion}
%% ============================================================

The MiniBooNE/LSND anomaly remains \emph{ambiguous} pending final SBN
results.  The RS analysis provides one firm structural conclusion: there
is no fourth neutrino mass eigenstate.  The dimension $D = 3$ of the
recognition framework forces exactly three active neutrino generations,
and a sterile neutrino at $m_4 \sim 1$~eV is structurally excluded.
The $\varphi$-aliasing hypothesis offers a qualitative mechanism for
residual excess events, but quantitative derivation remains as future
work.  MicroBooNE data supporting background explanations are consistent
with the RS prediction.  The SBN program will provide the definitive
test.

%% ============================================================
\begin{thebibliography}{99}

\bibitem{RS_Axioms}
J.~Washburn, ``The Algebra of Reality: A Recognition Science Derivation
of Physical Law,'' \emph{Axioms} \textbf{15}(2), 90 (2026).
\texttt{doi:10.3390/axioms15020090}

\bibitem{Aczel1966}
J.~Acz\'el, \emph{Lectures on Functional Equations and Their
Applications} (Academic Press, 1966).

\bibitem{PDG}
S.~Navas \textit{et al.}\ (Particle Data Group),
``Review of Particle Physics,'' Phys.\ Rev.\ D \textbf{110}, 030001 (2024).

\bibitem{LSND}
A.~Aguilar et al.\ (LSND Collaboration),
``Evidence for neutrino oscillations from the observation of
$\bar\nu_e$ appearance in a $\bar\nu_\mu$ beam,''
Phys.\ Rev.\ D \textbf{64}, 112007 (2001).

\bibitem{MiniBooNE2020}
A.~A.~Aguilar-Arevalo et al.\ (MiniBooNE Collaboration),
``Updated MiniBooNE neutrino oscillation results with increased data
and new background studies,''
Phys.\ Rev.\ D \textbf{103}, 052002 (2021).

\bibitem{MicroBooNE2021}
P.~Abratenko et al.\ (MicroBooNE Collaboration),
``Search for an anomalous excess of inclusive charged-current
$\nu_e$ interactions in the MicroBooNE experiment using Wire-Cell
reconstruction,''
Phys.\ Rev.\ D \textbf{105}, 112005 (2022).

\end{thebibliography}

\end{document}
